% Sources: prism/report/facts_v3_1.py; prism/report/prompts_v3_1.py; prism/report/checker.py;
% prism/report/verifier.py; REPORT_V3.md; REPORT_V3_1.md; METHODS_RESULTS.md M3; CHECKER_MUTATION.md;
% REVISION_DATA.md R3/R4; refs kwon2023, herasrivera2026, weinberg2018, gore2019, zeba2026, wiseman2017, gabriel2021,
% rao2025rexerr.
\paragraph*{Record, report and checks.} PRISM writes everything it computes for a scan into one structured record,
and the report is written from that record alone. The record is a flat list of named entries. Each entry has a key,
such as \code{meas.enhancing\_cc}, and holds one value or one fixed sentence. Together they give the time after
radiotherapy, the volumes and their change since the previous scan, each model's score with its cohort median and the
inputs it is based on, the conformal sets, the agreement statement, the reliability card and four flags. A flag is
raised when the clinical-history and imaging models disagree, when IDH or MGMT status is unknown or an input lies
outside the 1st to 99th percentile of its training values, when the card has few cases or none for the scan's group, or when the combined score's 80\% set
holds both outcomes. Numbers are rounded as the record is built (scores to two decimals, volumes to one). Every wording
that can be fixed in advance, such as the tier, each conformal set, the agreement, each change in volume and the
reason for each flag, is stored as a sentence for the language model to copy. The record holds no outcome label.

A local open-weight language model, Qwen3.8-27B, writes the report's prose from the record. It runs offline on one GPU
of the computer that runs PRISM, served with the vLLM library \citep{kwon2023}, so the record never leaves that
computer. Decoding is greedy: the model always takes its most likely next word, so the report does not depend on
random sampling. The prompt asks for sentences in a fixed order: the scan's timing; the volumes, called indicative;
their change; each model's score and what it is based on; whether the two models agree; the combined score and what a
score means; its 80\% conformal set; the card's tier; and one sentence for each raised flag. It tells the model to copy
every number as digits and every fixed phrase exactly, and to add no diagnosis, response category, advice or
probability words such as ``likely''. The model sees no images and has no tools, and every number it may write is
already in the record, so it has nothing to compute. Its output is constrained to a JSON schema: a list of sentences,
each with its text, one of four report sections and the keys of the record entries it uses. The schema forces this
shape, so every report can be split into sentences with their cited keys. It does not force a cited key to exist or
the text to agree with it, which is left to the checks below.

Fixed code, not the language model, fills in the rest of the report page: the tables of volumes and scores, the
conformal sets, the card, the flags, the segmentation slices and a fixed paragraph of limitations. The language model
only turns record entries into prose, a division also used to write brain tumour reports from extracted image
features \citep{herasrivera2026}. The page carries the history, comparison and change elements of the BT-RADS
follow-up template \citep{weinberg2018}, whose structured reports were rated more consistent than free prose
\citep{gore2019}. It leaves out the template's category and advice, because PRISM makes no decision (Supplementary
Section~S8). Because the enhancing volume varies between scans taken close together (Section~\ref{sec:3.3}), the
report calls all volumes indicative.

Two rule-based programs then compare the report with its record, and neither uses a language model. The first, the
checker, tests each sentence against the entries it cites and the report as a whole against the record. It rejects a
report that cannot be parsed, cites a key that is not in the record, writes a number that matches no cited entry up to
the rounding written, or writes a number as a word. It also rejects a report that leaves out a raised flag or calls an
input out of range when no flag says so, misstates a conformal set or whether the models agree, leaves out any of the
three scores or the enhancing volume, omits the tier where the card has few cases or none, quotes a performance figure
that is not on the card, or uses words of diagnosis, response category, advice, probability or certainty. Rule-based
checks of this kind test generated text against the source record \citep{zeba2026}, as extractive evaluation of text
generated from tables does with a learned extractor \citep{wiseman2017}. The checker was fixed before any evaluation
report was generated, and a rejected report is not edited or written again.

The prompt was revised once. In the first evaluation the checker rejected all 178 LUMIERE reports because of one
sentence: the prompt asked for the rater's three-month mark in words that contain the number 3, but the record entry
held only ``yes'' or ``no'', so the 3 was in no cited entry. We then changed the record so that this entry holds the
full sentence with its number and every number word is written as digits, made each line of the prompt appear only
when its entry is in the record, and removed the mention of RANO. The revised prompt was developed on 58 reports for
20 MU, 10 LUMIERE and 10 UCSF-ALPTDG patients, and the evaluation reports were then generated again for the other
patients with the checker unchanged.

The checker accepts a number that equals a number in any cited entry, and it does not read direction words. It
therefore accepts a sentence that gives the FLAIR volume as the enhancing volume while citing both entries, or that
writes ``larger'' where the record says ``smaller''. The second program, the verifier, runs the checker and adds two
rules that tie each number and direction word to what it describes. It splits each sentence into clauses at commas,
semicolons, colons and the word ``and'', and assigns each number and direction word to the nearest tumour region or
score named in its clause. The first rule requires each such number to equal, up to its written rounding, a value that
the cited entries give for that region or score. The second requires each direction word on a fixed list (larger,
increased, grew, rose, higher, and their opposites) to agree with the direction that the cited entries give for that
region. As first written, the rules took an entry's region from its key alone. They rejected every MU and LUMIERE
evaluation report and 24 of 52 correct synthetic reports, because most entries name their region in their text, as in
``enhancing tumour 3.5\,cc larger than at the previous scan''. We let the rules also read the region from the text of a
cited entry, tested this change only on the synthetic records and the 48 MU and LUMIERE development reports, fixed it,
and then ran it once on the 276 MU and LUMIERE evaluation reports.

To measure what the two programs catch, we planted one error at a time in correct reports, with 24 kinds of error
(Supplementary Section~S10), as planted errors are used to evaluate factuality checks \citep{gabriel2021,rao2025rexerr}.
The correct reports came from three writers: test code, for 52 synthetic records that cover two tiers, both sides of
each model's median, each of the three possible conformal sets and the out-of-range flag; the language model, for the
276 MU and LUMIERE evaluation records; and a fixed template that fills one set sentence for each line of the prompt
from the same records. Three kinds of error, two values placed under the wrong region or score and one change
rewritten in the opposite direction, were held out while the rules were written. Because we wrote both the rules and
these errors, a separate language model that had not seen the rules then rewrote the synthetic reports with errors in
its own words, with correct paraphrases as controls (exploratory). To test whether a report follows its record, we
changed one entry in 30 evaluation records, six each for the card's tier, the agreement, the change in enhancing
volume, the out-of-range flag and the conformal set, and planned that at least 27 reports would state the new value
and not the old one. We also gave the same prompt, unchanged, to six other local open-weight models of 2 to 35 billion
parameters.

Both programs test whether a report agrees with its record. Neither tests whether the record, or the report, is
clinically correct: that rests on the scores and measurements, whose evaluation follows.
